% Folder 77, batch 08 — archivists' pages 141–157.
% Pass: Opus 5.5 (claude-opus-5-5), 2026-10-03 — first pass, unchecked against the pages by a human.
%
% Notation: his underlined scheme-valued objects are set with underline in math;
% R_i = E_i/E_{i-1}, Q = E_nu^perp/E_nu; his double-barred mu_2 is mu_2.
\documentclass[11pt,a4paper]{article}
\input{../preamble/grothendieck}

\folder{77}
\batch{8}
\pages{141}{157}
\foldertitle{Quadriques affines, extensions quadratiques : notes manuscrites (s.d.).}
\dating{[à partir de 1977]}
\watermark{Édition de démonstration}

\begin{document}

\page{141}
\note{la page commence au milieu d'un raisonnement venu du lot précédent.}
ramené : prouver que ce dernier est connexe. Or il est fibré sur
$\underline{\Delta}_{\{n-1\}}$ (qui est connexe), il suffit donc de prouver

\emph{Lemme 6} $\underline{\Delta}^{\omega}_{\{n-1,n\}} \xrightarrow{\ \sim\ } \underline{\Delta}_{\{n-1\}}$
par la projection canonique, en d'autres termes
$\forall F'\in\Delta_{\{n-1\}}(E,q)$, $\exists$ exactement un
$F\in\Delta^{\omega}_{\{n\}}(E,q)$ tel que $F\subset F'$.
\marginal{indépendant de la question de car.}
[On connaît les fibres de $(\underline{\Delta}_{\{n-1,n\}})|F$ au-dessus de
$\underline{\mathrm{Arf}}$ \ill{} \ill{} \uncertain{iso}]
\note{le « $F\subset F'$ » est écrit tel quel ; vu les rangs ($n$ et $n-1$), on attendrait $F'\subset F$.}

\emph{Dém.} Par le \uncertain{lemme 3}, on sait que l'image inverse de $F'$ dans
$\underline{\Delta}_{\{n,n-1\}}$ est \uncertain{iso} à
$\underline{\Delta}_1(F'^{\perp}/F')$, donc est un rev. étale de degré $2$ de $S$.
\note{sous « $F'^{\perp}/F'$ » il a écrit « rg $2$ ».}
Il suffit de voir que les morphismes \add{de celui-ci} dans $\underline{\mathrm{Arf}}$
commutent aux automorphismes \uncertain{can.} de ces rev., ce qui est
\uncertain{évident}.

Une fois traité le cas de corps \uncertain{F}, \struck{on obtient}
\add{on déterminera} les ens. de comp. connexes de $\underline{\Delta}_{\{n\}}$, et on
\uncertain{passera} au cas général par un argument de spécialisation (fact. de
Stein …). En même temps, on définit $\underline{\mathrm{Arf}}(E,q)$ dans le cas
général, par factorisation de Stein de $\underline{\Delta}_{\{n\}}$, et on verra qu'il
\uncertain{marche} pour tous les $\underline{\Delta}_H$ avec $H\ni n$.

\page{142}
Récapitulons

\emph{Thm} $(E,q)$ forme quadratique régulière de rg
\struck{$E=\ill{}$}, $N=2n$ \struck{\ill{}} ou $N=2n+1$ ($n\geq 1$, donc $N\geq 2$).
Soient $H,H'\subset[1,n]$, $d=\max H$, $d'=\max H'$ (donc $0\le d,d'\le n$), considérons
\[
(1)\qquad \underline{\Delta}_H(E,q) \xrightarrow{\ p_{H',H}\ } \underline{\Delta}_{H'}(E,q)
\]
\note{on attend $H'\subset H$ pour que la projection (1) soit définie ; l'inclusion n'est pas écrite ici.}

a) Ce morphisme est projectif, \struck{et} lisse, surjectif.
En particulier (prenant $H'=\emptyset$) les $\underline{\Delta}_H(E,q)\to S$ sont
projectifs, lisses, surjectifs.

b) Les fibres géométriques de ce morphisme sont \struck{\ill{}} connexes, sauf dans le
cas $N$ pair, $d'<d=n$, auquel cas les fibres géométriques \struck{ont} ont deux
comp. conn.

c) Dans le \struck{dernier cas} \add{cas $N=2n$}, considérons
\[
\underline{\mathrm{Arf}}(E,q) = \text{schéma des comp. conn. de } \underline{\Delta}_{\{n\}}(E,q) \text{ sur } S
\]
qui est un rev. étale de \struck{degré} $2$ de $S$.
Considérons \add{(si $H\ni\{n\}$)} le \struck{facteur} morphisme composé
\[
\underline{\Delta}_H \longrightarrow \underline{\Delta}_{\{n\}} \longrightarrow \underline{\mathrm{Arf}}
\]
\struck{C'est lui qui ratisfait} Ce morphisme est
\page{143}
(projectif lisse) $0$-connexe (donc fournit la factorisation de Stein de
$\underline{\Delta}_H\to S$).

Plus généralement, si $H'\subset H$, $d'<d$, le morphisme
\[
(2)\qquad \underline{\Delta}_H \longrightarrow \underline{\Delta}_{H'}\times_S \underline{\mathrm{Arf}}
\]
est (projectif lisse) $0$-connexe (donc fournit, cf. (1), la factorisation de
Stein de $(H;H')$).\note{lecture de la fin de la phrase douteuse : « (donc fournit cf. … de $(H;H')$) (1) » ; l'ordre exact des mots n'est pas sûr.}

d) Supposons toujours $N=2n$, le morphisme
\[
(3)\qquad \underline{\Delta}_{\{n,n-1\}} \longrightarrow \underline{\Delta}_{\{n-1\}}
\]
est étale fini de degré $2$, plus précisément le \struck{factorisation de}
\uncertain{simplexe} (1) donne un iso
\[
\underline{\Delta}_{\{n,n-1\}} \xrightarrow{\ \sim\ } \underline{\Delta}_{\{n-1\}}\times \underline{\mathrm{Arf}}
\qquad (\text{pour } N=2n).
\]

e) Soit \struck{\ill{}}
\[
(4)\qquad \underline{\mathrm{Arf}}' = \underline{M}\bigl((-1)^n\delta(q)^{-1}\bigr),
\]
le \uncertain{sous-schéma} de $\underline{\det} E$ défini par
$\omega^{\otimes 2}=(-1)^n\delta(q)^{-1}$. Je \struck{trouve} \add{existe} un
morphisme unique
\[
(5)\qquad \underline{\mathrm{Arf}} \longrightarrow \underline{\mathrm{Arf}}'
\]
tel que le composé \uncertain{avec} $\underline{\Delta}_{\{n\}}\to\underline{\mathrm{Arf}}$ soit
\[
(6)\qquad
\begin{cases}
\underline{\Delta}_{\{n\}} \longrightarrow \underline{\mathrm{Arf}}' \\
F \longmapsto \omega_F \in \underline{\det}(E)^{*} \simeq \underline{\mathrm{Hom}}\bigl(\underbrace{\underline{\det} F\otimes\underline{\det}\check F}_{\underline{O}_S},\ \underline{\det} E\bigr)
\end{cases}
\]

\page{144}
déduit de la structure d'extension. Cet homomorphisme (5) est compatible aux
\uncertain{opérations} de $\mathbb{Z}/2\mathbb{Z}$, et définit donc un iso
\[
(7)\qquad \underline{\mathrm{Arf}}' \simeq \underline{\mathrm{Arf}} \overset{(\mathbb{Z}/2)_S}{\wedge} \mu_{2,S}
\]
\uncertain{En particulier}, si $2$ inversible sur $S$, on a un iso
\[
\underline{\mathrm{Arf}} \xrightarrow{\ \sim\ } \underline{\mathrm{Arf}}' = \underline{M}\bigl((-1)^n\delta(q^{*})^{-1}\bigr)
\]
\struck{f) \ill{} \ill{} \ill{} \ill{} \ill{} $\underline{O}(q)$}

\section*{Application aux groupes orthogonaux}

\textbf{Théorème} a) Supposons $N=2n+1$. Alors
\[
\underline{O}(q) \simeq \mu_2\times\underline{SO}(q) \qquad (\mu_2 = \text{\uncertain{centre}})
\]
et $\underline{SO}(q)$ est lisse à fibres connexes.
\struck{b) $\underline{SO}(q)$} $\underline{SO}(q)$ opère sur chaque
$\underline{\Delta}_H(E,q)$ via $\underline{SO}(q)\simeq\underline{O}(q)/\mu_2$,
\struck{et stabilise} un sp. homogène.\marginal{$\to \underline{SO}(q)$, le stabilisateur d'un \uncertain{drapeau} est \uncertain{un sous-groupe} lisse à fibres connexes.}
La dimension relative \ill{} \ill{} \ill{} \uncertain{produit} \ill{} :
\ill{} \ill{} \ill{} \ill{} \uncertain{en cat.} \uncertain{par ailleurs},
\ill{} \ill{} \struck{\ill{}} ;

\uncertain{Cor.} La \uncertain{projection} $\underline{O}(q)\to\mu_2$ est
\uncertain{induite par} $\det : \underline{O}(q)\to\mathbb{G}_m$.

\struck{\ill{} \ill{} \ill{} \ill{} \ill{} \ill{} \ill{} $\underline{SO}$, $\underline{O}(q)$ qui \ill{} \ill{} $F\in\Delta_{\{n\}}(E,q)$.}

\page{145}
\note{la page entière est barrée de grands traits obliques ; on la transcrit telle qu'elle se lit.}
\struck{Alors $SO(q)$ est fibré sur}
[Considérons
\[
\begin{cases}
\underline{SO}(q) \longrightarrow \underline{\Delta}_{\{n\}} \\
u \longmapsto u(F)
\end{cases}
\]
il s'agit de prouver que ce morphisme est lisse à fibres géom. connexes. Plus
généralement] \struck{\uncertain{Unicité}}

[$D\in\underline{\Delta}_H(E)$
\[
\begin{cases}
\underline{SO}(q) \longrightarrow \underline{\Delta}_H \\
u \longmapsto u(D)
\end{cases}
\]
est lisse à fibres géom. connexes (dans $\underline{\Delta}_H$ \uncertain{espace}
homogène sous $\underline{SO}(q)$)]. Par transitivité OPS $H=[1,n]$.

\uncertain{Comme} $\underline{SO}(q)$ \struck{\ill{}} \uncertain{opère transitivement}
sur $\underline{\Delta}_H$, $\underline{O}(q)$ \ill{} \ill{}.
\struck{\ill{}}\; $D=\{E_1\subset\cdots\subset E_n\subset E\}$ \ill{} $\xi=1$.
(cf. \ill{} \ill{} \ill{})

On veut ainsi \ill{} (\ill{} \ill{}) \ill{} : une base standard
$e_0, e_1,\ldots,e_n, e_1^{*},\ldots,e_n^{*}$, \struck{$e_0$}
$E_i$ engendré par $e_1,\ldots,e_i$. \struck{A\ill{}}
Si $D'=(E'_1\subset\cdots\subset E'_n)\in\Delta_{[1,n]}$ \ill{}
$e'_0, e'_1,\ldots,e'_n, e'^{*}_1,\ldots,e'^{*}_n$.
Il y a un isomorphisme (unique) de $(E,q)$
\[
e_i\mapsto e'_i,\qquad e_i^{*}\mapsto e'^{*}_i ,
\]
\struck{Il donc}

\page{146}
\textbf{Proposition} $(E,q)$ forme quadratique régulière [$\varphi$ forme bil.
associée], sur un schéma $S$, $F$ sous-\struck{fibré} \add{(Module strict)}
\struck{tot.} isotrope de $E$.

a) L'hom. composé
\[
E \xrightarrow[\text{défini par }\varphi]{\ \psi\ } E^{\vee} \xrightarrow{\ \text{restr}\ } F^{\vee}
\]
est épi, donc son noyau $F^{\perp}$ est un sous-\struck{fibré} Module strict de
$E$, dont la formation commute aux changements de base.

b) [Soit $G=E/F^{\perp}$] La forme $\varphi$ induit un accouplement
\[
F\times G \longrightarrow \underline{O}_S
\]
qui est une dualité, et induit des isom.
\[
G \xrightarrow{\ \sim\ } \check F,\qquad F \xrightarrow{\ \sim\ } \check G .
\]

\note{à partir d'ici, la suite de la page (c)) est barrée de deux grands traits croisés.}
c) Soit $\underline{\mathrm{Split}}(E,F^{\perp})$ \add{ou $\underline{\mathrm{Scind}}$,}
(le faisceau des scindages de la suite exacte
\[
0\longrightarrow F^{\perp} \longrightarrow E \longrightarrow G \longrightarrow 0,
\qquad F\subset F^{\perp}
\]
qui est un torseur sous $\underline{\mathrm{Hom}}(G,F^{\perp})$, \add{ou $\underline{\mathrm{Splitisot}}$,}
soit $\underline{\mathrm{Splitisot}}(E,F^{\perp})$ le sous-faisceau des scindages
\struck{\ill{}} isotropes, i.e. tels que $\theta(G)$ soit \struck{tot.} isotrope dans
$E$, \struck{\ill{}}
\[
\underline{\mathrm{Splitisot}}(E,F^{\perp}) \hookrightarrow \underline{\mathrm{Split}}(E,F^{\perp}).
\]
Considérons d'autre part la suite exacte
\[
0\longrightarrow \underline{\mathrm{Hom}}(G,F) \longrightarrow \underline{\mathrm{Hom}}(G,F^{\perp}) \longrightarrow \underline{\mathrm{Hom}}(G,Q) \longrightarrow 0
\]
où
\[
Q = F^{\perp}/F .
\]

\page{147}
\struck{c) Soit} \add{d) Soit} $Q=F^{\perp}/F$, alors $\varphi$ induit \uncertain{une forme
quadratique} $q_Q$ sur $Q$, et \struck{la forme bil. associée} \add{celle-ci est
régulière}. \struck{Soit \ill{}} \add{forme}
\[
\theta : G \longrightarrow E
\]
un scindage de la suite exacte
\[
(*)\qquad 0\longrightarrow F^{\perp} \longrightarrow E \longrightarrow G \longrightarrow 0
\]
et considérons la forme quadratique
\[
q_\theta : F\times G \longrightarrow \underline{O}_S
\]
définie par
\[
q_\theta(x,y) = q\bigl(x+\theta(y)\bigr) = q(\theta(y)) + \langle y,x\rangle .
\]
La forme bilinéaire associée est
\[
\begin{aligned}
\varphi_\theta(x,y;x',y') &= \varphi(x+\theta y,\theta y') + \varphi(x',\theta y)\\
&= \varphi(\theta y,\theta y') + \langle y',x\rangle + \langle y,x'\rangle
\end{aligned}
\]
et elle est non \emph{dégénérée}, donc $q_\theta$ est une
\marginal{et \struck{la forme} $q|Q_\theta$ est régulière, [\emph{NB} $Q_\theta\simeq Q=F^{\perp}/F$, et la forme $q_Q$ déduite de $q_\theta$ ne dépend pas du choix de $\theta$ — \uncertain{c'est trivial} : \ill{} \ill{} \ill{} …]}
forme quadratique régulière (de rg pair). Si $Q_\theta$
est l'orthogonal de $F+\theta(G)$, on a donc
$E \simeq (F\oplus\theta(G))\oplus Q_\theta$.

d) Supposons $S$ affine, \struck{\uncertain{ce qui}} \uncertain{garantit} l'existence
d'un $\theta$…. Alors il existe (pour $\theta$ fixé)
\[
u : G \longrightarrow F
\]
tel que, posant
\[
\theta_u = \theta + u
\]
on ait
\[
q\circ\theta_u = 0
\]
\marginal{\emph{NB} condition qui équivaut à : $(q\circ\theta)(y)+\langle y,u(y)\rangle = 0$, i.e. $B_u(y,y) = -(q\circ\theta)(y)$, où $B_u$ est la forme bil. associée à $u$.}
i.e. $\theta_u(G)$ soit tot. isotrope dans $E$ (scindage isotrope). L'ensemble des
\struck{$u\in\underline{\mathrm{Hom}}($}
\[
u \in \mathrm{Hom}(G,F) \simeq \mathrm{Bil}(G,G) \qquad (F\simeq\check G)
\]
tels que \struck{$q\circ$} $\theta_u$ soit isotrope, est un torseur
sous le sous-groupe $\mathrm{Alt}(G,G)$ de $\mathrm{Bil}(G,G)$.
\struck{Il\ill{}} \uncertain{Indépendamment} du choix d'un $\theta$)

\page{148}
l'ensemble $\mathrm{Scind\,isot}(E,F^{\perp})$ des scindages isotropes de $(*)$ est
stable par les opérations de
\[
\mathrm{Alt}(G,G) \hookrightarrow \mathrm{Bil}(G,G) \simeq \mathrm{Hom}(G,F) \hookrightarrow \mathrm{Hom}(G,F^{\perp})
\]
et le quotient $\mathrm{Scind\,isot}/\mathrm{Alt}(G,G)$ est isomorphe (par la flèche
canonique) à $\mathrm{Scind}/\underbrace{\mathrm{Bil}(G,G)}_{\mathrm{Hom}(G,F)}$ (qui est lui-même un
torseur sous $\mathrm{Hom}(G,Q)$), où $Q=F^{\perp}/F$ …

e) Soit $\theta$ un scindage isotrope de $(*)$. Alors $(E,q)$ est somme directe
quadratique de $F+\underbrace{\theta(G)}_{\check F}$, \struck{décomposé} où $q$ est somme
forme décomposée
\[
q(x,x') = \langle x,x'\rangle ,
\]
et de $Q_\theta$.

\textbf{Corollaire} Considérons un drapeau
\[
D :\ 0 \subsetneq E_1 \subsetneq E_2 \subset\cdots\subsetneq E_\nu = F \qquad E ,
\]
d'où un drapeau complété
\[
0\subsetneq E_1 \subset\cdots\subsetneq E_\nu = F \subset F^{\perp} \subsetneq E_{\nu-1}^{\perp} \subsetneq\cdots\subsetneq E_1^{\perp} \subsetneq F
\]
\note{le dernier terme se lit $F$ sur la page, là où l'on attend $E$.}
[\emph{NB} on peut avoir $F=F^{\perp}$ donc $Q=F^{\perp}/F=0$ …]. Un
scindage isotrope de $(*)$ définit alors une structure \uncertain{scindée} dans
$E$ identifié à : $F\oplus\check F\oplus Q$.

\page{149}
\textbf{Corollaire 1} Soient $(E,q)$, $(E',q')$ formes quadr. rég. de
\uncertain{même type} [$N=2n$ ou $N=2n+1$] sur $S$, \struck{$F\subset E$} et
\[
D\in\Delta_H(E,q),\qquad D'\in\Delta_H(E',q') ,
\]
$H\subset[1,n]$, enfin $D=\{E_1,\ldots,E_\nu\}$, $D'=\{E'_1,\ldots,E'_\nu\}$.
\struck{Supposons, et supposons} Posons $Q=E_\nu^{\perp}/E_\nu$,
$Q'=E'^{\perp}_\nu/E'_\nu$, et soient $q_\Theta$ et $q_{\Theta'}$
\ill{} \ill{} \ill{} \ill{} sur $S$. Si $q_\Theta\simeq q_{\Theta'}$,
\add{et si $R_i=E_i/E_{i-1}\simeq R'_i=E'_i/E'_{i-1}$,} et si $S$ est affine, alors
\struck{\ill{}}
\[
\exists\, u : (E,q) \xrightarrow{\ \sim\ } (E',q') \quad\text{tel que}\quad u(D)=D' .
\]
\uncertain{On en conclut}

\textbf{Corollaire 2} Il existe loc. (ét) un isom.
\[
u : (E,q) \longrightarrow (E',q') \quad\text{tel que}\quad u(D)=D' .
\]

Donc,

\textbf{Corollaire 3} $\underline{\Delta}_H(E,q)$ est un esp. homogène
\uncertain{(ét)} sous $\underline{\mathrm{Aut}}(E,q)=\underline{O}(q)$. En d'autres
termes, si $D\in\Delta_H(E,q)$, alors
\[
\begin{cases}
\underline{O}(q) \longrightarrow \underline{\Delta}_H \\
u \longmapsto u(D)
\end{cases}
\]
est un épim. (ét), \struck{identification} définissant un iso
\[
\underline{\Delta}_H \xleftarrow{\ \sim\ } \underline{O}(q)/\underline{\mathrm{Stab}}_D(q)
\]
où $\underline{\mathrm{Stab}}_D(q)$ est le stabilisateur de $D$ dans $\underline{O}(q)$.

\page{150}
\emph{Corollaire}

Nous pouvons préciser le cor. 1 ainsi :

\emph{Cor. 4} Sous les conditions du cor. 1, posons
\struck{$R_i$ \quad $1\le i\le\nu$ donnons \ill{} pour \quad $R_i$ \quad $1\le i\le\nu$ \ill{} \ill{}}
\[
R_i = E_i/E_{i-1},\qquad R'_i = E'_i/E'_{i-1}\qquad (E_0=0,\ E'_0=0)
\]
et donnons-nous les isom.
\[
\begin{aligned}
u_i &: R_i \xrightarrow{\ \sim\ } R'_i\\
u_q &: Q \xrightarrow{\ \sim\ } Q'
\end{aligned}
\]
\add{(ce dernier)} orthogonal pour $q$, $q'$. Alors, \struck{il existe} \add{si $S$ est affine,}
a) il existe \add{(ici iso.)}
\[
u : (E,q) \longrightarrow (E',q')
\]
tel que $u(D)=D'$, et que $u$ induise les $u_i$ et $u_q$.

\struck{\uncertain{b)}} L'ens. de ces $u$ est \add{(évident)} un torseur à dr.
sous le \uncertain{sous-}groupe $\underline{O}_D(q)^{+}$ des
\struck{$\to\ \underline{O}_D(q)$} \add{$\subset$} $\underline{\mathrm{Aut}}(E,q,D)$
$= \underline{O}_D(q)(S)$, formé des automorphismes qui induisent l'identité sur
les $G_i$ et sur $Q$.

\struck{Considérons d'autre} \add{Considérons d'autre part} l'ens. \struck{des} produit
\[
\overline{\Sigma} = \prod_{1\le i\le\nu} \Sigma_i \times \Sigma_0
\]
où $\Sigma_i$ est l'ens. des scindages de
\[
0\longrightarrow E_{i-1} \longrightarrow E_i \longrightarrow R_i \longrightarrow 0
\]
et $\Sigma_0$ l'ens. des scindages \emph{isotropes} de
\[
0\longrightarrow E_\nu \longrightarrow E_\nu^{\perp} \longrightarrow Q \longrightarrow 0 .
\]
\struck{Considérons d'autre part} \add{\ill{}}
\[
\overline{E} = \bigoplus_i \bigl(R_i + \check R_i\bigr) \oplus Q
\]
avec la forme $\overline{q}$ somme des formes

\page{151}
évidentes $\langle x_i,$ \struck{$x_i$} $x'_i\rangle$ sur les $G_i\oplus\check G_i$, et $q_Q$
sur $Q$. \struck{Alors} Soit $\overline{D}=\{\overline{E}_1,\ldots,\overline{E}_\nu\}$,
avec
\[
\overline{E}_i = \sum_{1\le j\le i} R_j \subset \overline{E} .
\]
\struck{Considérons} Considérons
\[
\begin{cases}
\overline{R}_i = \overline{E}_i/\overline{E}_{i-1}\\
\overline{Q} = \overline{E}_\nu^{\perp}/\overline{E}_\nu
\end{cases}
\]
\uncertain{on a} donc des isom. canoniques
\[
\begin{aligned}
\overline{R}_i &\xrightarrow[\ \overline{u}_i\ ]{\ \sim\ } R_i\\
\overline{Q} &\xrightarrow[\ \overline{u}_q\ ]{\ \sim\ } Q ,
\end{aligned}
\]
\uncertain{lesquelles} \uncertain{étant} compatibles avec $\overline{q}$, $q$.

Ceci posé … :

\emph{Cor. 5} L'ens. $\overline{\Sigma}$ est canoniquement isomorphe à
l'ens. des isom. $(\overline{E},\overline{q},\overline{D}) \to (E,q,D)$
compatibles avec les $\overline{u}_i$, $\overline{u}$.

En termes schématiques, le schéma
produit $\underline{\overline{\Sigma}} = \prod_{1\le i\le\nu}\underline{\Sigma}_i\times\underline{\Sigma}_0$
est canoniquement \marginal{\uncertain{sauf} \uncertain{vérif.} $S$ affine !}
isomorphe au schéma
\[
\underline{\mathrm{Isom}}_{\{\overline{u}_i\},\overline{u}}\bigl((\overline{E},\overline{q},\overline{D});(E,q,D)\bigr)
\]
Ce dernier est un torseur : \uncertain{à gauche}
sous $\underline{O}^{+}_D(q)$ (\uncertain{à droite} sous $\underline{O}^{+}_{\overline{D}}(\overline{q})$).
Comme $\overline{\Sigma}$ est évidemment lisse et connexe
sur $S$, \uncertain{on} en conclut

\page{152}
\emph{Cor.} 6 Le \add{schéma en} groupes $\underline{O}_D(q)^{+}$ est lisse à
fibres connexes, plus précisément, \struck{si} \add{si} $S$ est affine
(\uncertain{car} $\overline{\Sigma}$ a une section donc $\underline{O}_D(q)^{+}\simeq\overline{\Sigma}$) il est
isom. comme schéma (mais sans structure de groupe) à un fibré \struck{affine} vectoriel
\[
\check{\mathbb{V}}\Bigl(\sum_{1\le i\le\nu}\underline{\mathrm{Hom}}(R_i,E_{i-1})\oplus\underline{\mathrm{Hom}}(\overline{G},\overline{F}^{\perp})^{!}\Bigr)
\]
\note{devant le second terme, un signe somme a été barré.}
où $\underline{\mathrm{Hom}}(\overline{G},\overline{F}^{\perp})^{!}$ est le noyau de l'épi.
\[
\begin{cases}
\underline{\mathrm{Hom}}(\overline{G},\overline{F}^{\perp}) \longrightarrow \underline{\mathrm{Quad}}(\overline{G})\\
u \longmapsto \bigl(\overline{y}\mapsto \overline{q}(\overline{y}+u(\overline{y}))\bigr)
\end{cases}
\]
\note{au-dessus de la seconde ligne, une première écriture du membre de droite, commençant par $\overline{q}$ et se terminant par $+\varphi($, est raturée.}

Calculons les dimensions, posant
\[
\operatorname{rg} R_i = d_i,\qquad \operatorname{rg} Q = \rho \qquad\text{donc}\quad N = 2\overbrace{\textstyle\sum d_i}^{d}+\rho,
\]
\[
\operatorname{rg}\overline{F} = \operatorname{rg} G = \sum d_i \overset{\text{déf}}{=} d
\]
\[
\begin{cases}
\dim_S \underline{\Sigma}_i = d_i \displaystyle\sum_{1\le j\le i-1} d_j\\[1ex]
\dim_S \underline{\Sigma}_0 = \dim_S \underline{\mathrm{Hom}}(\overline{G},\overline{F}^{\perp}) - \dim_S \underline{\mathrm{Quad}}(\overline{G})\\[1ex]
\phantom{\dim_S \underline{\Sigma}_0} = d(d+\rho) - \dfrac{d(d+1)}{2} = \dfrac{d(d-1)}{2}+d\rho
\end{cases}
\]
\note{après le dernier « $=$ » de la seconde accolade, un premier début ($\sum$ …) est raturé.}
\[
\begin{aligned}
\dim_S \underline{O}_D(q)^{+} &= \underbrace{\sum_{1\le j<i\le\nu} d_i d_j}_{\frac12(d^2-\sum d_i^2)} + \frac{d(d-1)}{2} + d\rho\\
&= d^2 + d\rho - \sum\frac{d_i(d_i+1)}{2}
\end{aligned}
\]
\note{dans la dernière ligne, un facteur $\tfrac12$ devant le signe somme est barré.}
\[
\dim_S\Bigl(\prod \mathrm{Gl}(R_i)\times\underline{O}(q_Q)\Bigr) = \sum d_i^2 + \frac{\rho(\rho-1)}{2}
\]
\[
\dim_S \underline{O}_D(q) = d^2 + d\rho + \frac{\rho(\rho-1)}{2} + \sum_i \frac{d_i(d_i-1)}{2}
\]
\note{ici encore un facteur $\tfrac12$ devant le dernier signe somme est barré.}
\[
\dim_S \underline{\Delta}_D(q) = \underbrace{\dim_S \underline{O}(q)}_{\frac{N(N-1)}{2} = 2d^2+d(2\rho-1)+\frac{\rho(\rho-1)}{2}} - \dim_S \underline{O}_D(q) = \text{\struck{$\frac{N(N-1)}{2}$}}
\]

\page{153}
\[
\dim_S \underline{\Delta}_H(q) = d^2 + d\rho - d - \sum\frac{d_i(d_i-1)}{2}
\]

Par exemple
\[
\dim_S \underline{\Delta}_{\{d\}}(q) = \frac{d(d+1)}{2} + d(\rho-1) \qquad (\rho = N-2d)
\]
\note{le signe devant $d(\rho-1)$ est surchargé, et le « $-1$ » écrit par-dessus un autre signe ; la lecture « $+\,d(\rho-1)$ » est douteuse. La formule précédente donnerait, pour $\nu=1$, $\frac{d(d+1)}{2}+d(\rho-1)$, ce qui concorde.}
\struck{Ainsi}, si \uncertain{$q$} est de rg $N=2n$, et $d=\{n\}$ \add{(donc $\rho=0$)}
\uncertain{on trouve}
\[
\frac{n(n+1)}{2}\ \text{\struck{$-$}}\ - n = \frac{n(n-1)}{2}
\]


\note{un trait horizontal sépare la suite.}

Supposons $N=2n+1$ — donc \uncertain{on sait que}
\[
\underline{O}(q) = \underline{SO}(q)\times\mu_2 ,\qquad
\text{\struck{où}}\ 
\begin{cases}
\mu_2 = \underline{O}(q)\cap\mathbb{G}_m = Z(\underline{O}(q))\\
\underline{SO}(q) = \underline{O}(q)\cap\underline{Sl}(E)
\end{cases}
\]
\struck{\ill{} $\underline{SO}(q)$} \struck{donc} $\underline{SO}(q)$ \add{donc} $\underline{O}(q)$ opère sur les
$\underline{\Delta}_H(q)$ via \struck{$\mu_2$} — $\mu_2$ opérant trivialement dans
$\underline{\Delta}_H$ — de sorte que l'on peut \uncertain{aussi} \uncertain{dire} que
les $\underline{\Delta}_H$ sont aussi des esp. homogènes sous $\underline{SO}(q)$.
\struck{Pour} \struck{\ill{}} Pour $\forall D\in\Delta_H$, le stabilisateur
\add{est $\underline{SO}_D(q)$} dans $\underline{SO}(q)$, \struck{donc}. On a donc
\[
\underline{SO}_D(q) = \underline{O}_D(q)\cap\underline{Sl}(E)
\]
Considérons l'hom. canonique
\[
\underline{O}_D(q) \xrightarrow{\ \alpha\ } \prod_{1\le i\le\nu}\underline{Gl}(R_i)\times\underline{O}(q_Q)
\]
\note{des traits obliques fins traversent le facteur $\prod\underline{Gl}(R_i)$, sans qu'on puisse dire s'il est biffé.}
(où, comme dessus, $Q = E_\nu^{\perp}/E_\nu$,
$D=\{E_1\subset\cdots\subset E_\nu\}$). On constate que

\page{154}
l'iso commutatif dans

\begin{tikzcd}[column sep=large]
\underline{O}_D(q) \arrow[r] \arrow[rr, bend right=30, "\det_E"'] & \underline{O}(q_Q) \arrow[r, "\det_Q"] & \mathbb{G}_m
\end{tikzcd}

donc
\[
\underline{SO}_D(q) = \alpha^{-1}\bigl(\underline{SO}(q_Q)\bigr) .
\]
On trouve donc

\emph{Proposition} \struck{$\underline{SO}_D(q)$ est isom.} \add{On a un hom. can.}

\struck{On a $\underline{SO}_D(q) \xrightarrow{\ \alpha\ }$} \struck{$\prod_{1\le i\le\nu}\underline{Gl}(R_i)\times{}$} \struck{$\underline{SO}(q_Q)$}
\struck{induit un}

\begin{tikzcd}[column sep=small, row sep=normal, nodes={font=\small}]
1 \arrow[r] & \underline{O}_D(q)^{+} \arrow[r] & \underline{O}_D(q) \arrow[r] & \prod_{1\le i\le\nu}\underline{Gl}(R_i)\times\underline{O}(q_Q) \arrow[r] & 1\\
1 \arrow[r] & \underline{SO}_D(q)^{+} \arrow[r] \arrow[u, "\wr"] & \underline{SO}_D(q) \arrow[r] \arrow[u, hook] & \prod_{1\le i\le\nu}\underline{Gl}(R_i)\times\underline{SO}(q_Q) \arrow[r] \arrow[u, hook] & 1
\end{tikzcd}

donc $\underline{SO}_D(q)$ est une extension de
$\prod\underline{Gl}(R_i)\times\underline{SO}(q_Q)$ par
$\underline{SO}_D(q)^{+}\simeq\underline{O}_D(q)^{+}$, \struck{d'où} \add{qui est}
\struck{lisse et à fibres connexes (comme \ill{} \ill{} ext. de $2$ groupes \ill{} \ill{} à fibres connexes \ill{} \ill{})}
\note{deux lignes ici sont barrées d'un trait horizontal et de hachures obliques.}

Supposons $E$ de rg pair $N=2n$. Alors,

\emph{Thm} ($\underline{SO}(q)$ ($=\underline{O}(q)\cap\underline{Sl}(E)$)) est un groupe
lisse à fibres connexes.

\emph{OPS} \uncertain{pour} démontrer qu'il existe $F\subset E$
\struck{\ill{}} isotrope de rg $n$, donc $F^{\perp}/F=Q$ de rg $1$,
\note{avec $N=2n$ et $\operatorname{rg}F=n$, on attendrait $Q=0$ ; « de rg $1$ » correspondrait au cas $N=2n+1$. La page porte bien « $N=2n$ » plus haut.}
donc $\underline{SO}(q_Q)=1$, donc \uncertain{(si $D=\{F\}$)}
$\underline{SO}_D(q)$ est extension de \struck{\ill{}} $\underline{Gl}(F)$ (qui est
connexe) par $\underline{SO}_D$ (\ill{}) \uncertain{lisse} à fibres connexes

\page{155}
donc est lisse à fibres connexes. \emph{D'autre part}
$\underline{SO}(q)/\underline{SO}_D(q)\simeq\underline{\Delta}_{\{n\}}(q)$ est lisse
à fibres géom. connexes, donc aussi $\underline{SO}(q)$.

\emph{Corollaire} Les $\underline{SO}_D(q)$ \add{($=B_D(q)$)} sont lisses à fibres
connexes. On a $\underline{O}_D(q)\simeq B_D(q)\times\mu_2$.

\marginal{\emph{Prop.} \struck{($S\neq\emptyset$)} Soit $H\subset[1,n]$, \add{$d=\sup H$,} alors
$\underline{SO}(q)$ est transitif sur $\underline{\Delta}_H(q)$ \struck{sauf si \ill{}} si $d\neq n$, et
\uncertain{aussi} si $S\neq\emptyset$. Si $d=n$, et si $\omega$ est une section de
$\underline{\mathrm{Arf}}(q)$ sur $S$, alors $\underline{SO}(q)$ est transitif sur
$\underline{\Delta}^{\omega}_H(q)$.}
\note{cette note marginale, écrite verticalement, est reliée par une accolade au paragraphe qui suit.}

Supposons enfin $N=2n$, de sorte que $\mu_2 = Z(\underline{O}(q))\subset\underline{Sl}(E)$.
On ne va pas définir $\underline{SO}(q)$ comme $\underline{O}(q)\cap\underline{Sl}(E)$, à
cause de car. $2$. Mais on fait opérer \ill{} \ill{} $\underline{O}(q)$ sur
$\underline{\mathrm{Arf}}(q)$, d'où un hom. can.
\[
\underline{O}(q) \xrightarrow{\ \varepsilon_q\ } (\mathbb{Z}/2)_S
\]
\marginal{explicitement : soit $\eta : (\mathbb{Z}/2)_S\to\mu_{2S}$ can., alors
$\eta\circ\varepsilon_q=\det_q$ (restriction de $\det_E$ à $\underline{O}(q)$).
\emph{Cor.} si $2$ inv. sur $S$, $\underline{SO}(q)=\underline{O}(q)\cap\underline{Sl}(E)$.}
et on pose
\[
\underline{SO}(q) = \operatorname{Ker}\varepsilon
\]
de sorte que la suite \uncertain{exacte}
\[
1\longrightarrow \underline{SO}(q) \longrightarrow \underline{O}(q) \xrightarrow{\ \varepsilon_q\ } (\mathbb{Z}/2)_S \longrightarrow 1
\]
Considérons comme dessus un $D$, d'où

\page{156}
$\underline{O}_D(q)$ et
\[
1\longrightarrow \underline{O}_D(q)^{+} \longrightarrow \underline{O}_D(q) \longrightarrow \prod_{1\le i\le\nu}\underline{Gl}(R_i)\times\underline{O}(q_Q) \longrightarrow 1
\]
On vérifie alors

\struck{\emph{Proposition} Posons $\underline{SO}_D(q)=\underline{O}_D(q)\cap\underline{SO}$, stabilisateur de $D$ dans $\underline{SO}(q)$. Alors $S$}
\note{ces deux lignes sont barrées d'un trait horizontal et d'un zigzag.}

\emph{Prop.} La restriction $\varepsilon_q\,|\,\underline{O}_D(q)$ est égale au
composé
\[
\underline{O}_D(q) \longrightarrow \underline{O}(q_Q) \xrightarrow{\ \varepsilon_{q_Q}\ } (\mathbb{Z}/2)_S ,
\]
donc $\underline{SO}_D(q)$ est l'image inverse de $\underline{SO}(q_Q)$ par
$\underline{O}_D(q)\to\underline{O}(q_Q)$, d'où \uncertain{un} \uncertain{diagramme} de
suites exactes
\struck{\emph{Cor.} $\underline{SO}_D(q)$ est une extension de}
\struck{\ill{}}

\begin{tikzcd}[column sep=small, nodes={font=\small}]
1 \arrow[r] & \underline{O}_D(q)^{+} \arrow[r] & \underline{O}_D(q) \arrow[r] & \prod_{1\le i\le\nu}\underline{Gl}(R_i)\times\underline{O}(q_Q) \arrow[r] & 1\\
1 \arrow[r] & \underline{SO}_D(q)^{+} \arrow[r] \arrow[u, "\wr"] & \underline{SO}_D(q) \arrow[r] \arrow[u, hook] & \prod_{1\le i\le\nu}\underline{Gl}(R_i)\times\underline{SO}(q_Q) \arrow[r] & 1
\end{tikzcd}

de sorte que $\underline{SO}_D(q)$ est une ext. de
$\bigl(\prod\underline{Gl}(R_i)\bigr)\times\underline{SO}(q_Q)$ par
$\underline{SO}_D(q)^{+}\simeq\underline{O}_D(q)^{+}$, qui est lisse à fibres conn.
\note{un trait vertical en zigzag dans la marge gauche accompagne ces deux lignes.}

\emph{Thm} $(E,q)$ forme quad. rég. de rg pair $2n$.
Alors $\underline{SO}(q)$ est lisse à fibres connexes.

OPS $\exists F\in\Delta_n(q)$, d'où section $\omega$ de $\underline{\mathrm{Arf}}(q)$ sur $S$.
Prenons ci-dessus $D=\{F\}$, \struck{\ill{} $F$ est de rg $n$} donc $Q=0$. \emph{Donc}
par ce qui précède, $\underline{SO}_D(q)$ est lisse à fibres connexes.

\page{157}
D'autre part $\underline{SO}(q)/\underline{SO}_D(q)\simeq\underline{\Delta}^{\omega}_{\{n\}}$ est
lisse à fibres connexes, donc aussi $\underline{SO}(q)$, cqfd.

\emph{Corollaire} Les $\underline{SO}_D(q)=B_D(q)$ sont lisses à fibres connexes.
\note{la page s'arrête ici ; le reste du feuillet est blanc. Le lot se termine sur ce corollaire.}
\end{document}
